\documentclass[11pt]{article}
\usepackage{amsmath,amsfonts,amssymb,amsthm,mathtools}
\usepackage[margin=2.5cm]{geometry}
\usepackage{hyperref}
\usepackage{enumitem}

\newcommand{\R}{\mathbb{R}}
\newcommand{\T}{\mathbb{T}}
\DeclareMathOperator{\Tr}{Tr}
\DeclareMathOperator{\divergence}{div}

\newtheorem{theorem}{Theorem}
\newtheorem{lemma}[theorem]{Lemma}
\newtheorem{proposition}[theorem]{Proposition}
\newtheorem{corollary}[theorem]{Corollary}
\newtheorem{definition}[theorem]{Definition}
\newtheorem{remark}[theorem]{Remark}

\title{Formal Statement and Decomposed Proof of the Global Steady State \\
of the Vlasov--Maxwell--Landau System}
\author{Reference document for formalization}
\date{}

\begin{document}
\maketitle

\tableofcontents
\newpage

%=============================================================================
\section{Definitions and Setup}
\label{sec:setup}
%=============================================================================

We work in dimension $d = 3$ for both position and velocity.

\begin{definition}[Periodic spatial domain]
\label{def:torus}
Let $\T^3 = (\R / \mathbb{Z})^3$ denote the three-dimensional flat torus. We write $|\T^3|$ for its volume (with the standard normalization $|\T^3| = 1$ or, more generally, an arbitrary positive constant depending on the period).
\end{definition}

\begin{definition}[Landau collision matrix]
\label{def:landau_matrix}
For $z \in \R^3 \setminus \{0\}$, define the \emph{projection matrix}
\[
\Pi(z) = I_3 - \frac{z z^\top}{|z|^2},
\]
where $I_3$ is the $3 \times 3$ identity matrix and $z z^\top$ is the outer product. The \emph{Landau collision matrix} (for Coulomb interactions in 3D) is
\[
A(z) = \frac{1}{|z|} \Pi(z) = \frac{1}{|z|}\left(I_3 - \frac{z z^\top}{|z|^2}\right).
\]
More generally, for an interaction potential characterized by a scalar function $\Psi(|z|) > 0$:
\[
A(z) = \Psi(|z|)\left(|z|^2 I_3 - z z^\top\right).
\]
\end{definition}

\begin{definition}[Landau collision operator]
\label{def:landau_operator}
For a distribution $f : \R^3 \to (0, \infty)$ that is sufficiently smooth and decays sufficiently rapidly as $|v| \to \infty$, the \emph{Landau collision operator} is
\[
Q(f, f)(v) = \nabla_v \cdot \int_{\R^3} A(v - w) \big[ f(w) \nabla_v f(v) - f(v) \nabla_w f(w) \big] \, dw.
\]
Equivalently, using $\nabla f = f \nabla \log f$, this can be written in \emph{score form}:
\[
Q(f, f)(v) = \nabla_v \cdot \left\{ f(v) \int_{\R^3} A(v - w) f(w) \big[ \nabla_v \log f(v) - \nabla_w \log f(w) \big] \, dw \right\}.
\]
\end{definition}

\begin{definition}[Vlasov--Maxwell--Landau system]
\label{def:vml}
Let $f(t, x, v) > 0$ be a distribution function with $(x, v) \in \T^3 \times \R^3$, and let $E(t, x), B(t, x) \in \R^3$ be the electric and magnetic fields. Let $\nu > 0$ be the collision frequency and $\rho_{\mathrm{ion}} > 0$ the uniform neutralizing ion background density. The \emph{Vlasov--Maxwell--Landau (VML) system} is:
\begin{align}
\label{eq:vlasov}
&\partial_t f + v \cdot \nabla_x f + (E + v \times B) \cdot \nabla_v f = \nu \, Q[f, f], \\
\label{eq:ampere}
&\partial_t E = \nabla_x \times B - J, \\
\label{eq:faraday}
&\partial_t B = -\nabla_x \times E, \\
\label{eq:gauss}
&\nabla_x \cdot E = \rho - \rho_{\mathrm{ion}}, \\
\label{eq:divB}
&\nabla_x \cdot B = 0,
\end{align}
where $J(t,x) = \int_{\R^3} v \, f(t, x, v) \, dv$ is the current density and $\rho(t,x) = \int_{\R^3} f(t, x, v) \, dv$ is the electron charge density.
\end{definition}

\begin{definition}[Local Maxwellian]
\label{def:local_maxwellian}
A \emph{local Maxwellian} is a distribution of the form
\[
f(x, v) = \frac{n(x)}{(2\pi T(x))^{3/2}} \exp\!\left(-\frac{|v - u(x)|^2}{2T(x)}\right),
\]
where $n(x) > 0$ is the number density, $u(x) \in \R^3$ is the bulk velocity, and $T(x) > 0$ is the temperature.
\end{definition}

\begin{definition}[Steady state]
\label{def:steady_state}
A \emph{steady state} of the VML system is a triple $(f, E, B)$ satisfying:
\begin{enumerate}[label=(\roman*)]
\item $\partial_t f = 0$, $\partial_t E = 0$, $\partial_t B = 0$;
\item $f(x, v) > 0$ for all $(x, v) \in \T^3 \times \R^3$;
\item $f$ is sufficiently smooth and, together with all its derivatives, decays sufficiently fast as $|v| \to \infty$;
\item the total energy $\mathcal{E} = \iint_{\T^3 \times \R^3} \frac{1}{2}|v|^2 f \, dv \, dx + \frac{1}{2}\int_{\T^3}(|E|^2 + |B|^2) \, dx$ is finite;
\item the entropy $\iint_{\T^3 \times \R^3} f \log f \, dv \, dx$ is finite.
\end{enumerate}
\end{definition}

\begin{definition}[Total energy]
\label{def:total_energy}
The \emph{total energy} of the VML system is
\[
\mathcal{E} = \underbrace{\iint_{\T^3 \times \R^3} \frac{1}{2}|v|^2 f \, dv \, dx}_{\text{kinetic}} + \underbrace{\frac{1}{2}\int_{\T^3}(|E|^2 + |B|^2) \, dx}_{\text{electromagnetic}}.
\]
\end{definition}

%=============================================================================
\section{Algebraic Lemmas about the Landau Matrix}
\label{sec:algebraic}
%=============================================================================

\begin{lemma}[Symmetry of $A$]
\label{lem:A_symmetric}
For all $z \in \R^3 \setminus \{0\}$:
\begin{enumerate}[label=(\alph*)]
\item $A(z)$ is a symmetric matrix: $A(z)^\top = A(z)$.
\item $A(z)$ is even: $A(-z) = A(z)$.
\end{enumerate}
\end{lemma}

\begin{proof}
(a) We have $\Pi(z) = I_3 - z z^\top / |z|^2$. Since $(z z^\top)^\top = z z^\top$, we get $\Pi(z)^\top = \Pi(z)$, so $A(z)^\top = A(z)$.

(b) Since $|-z| = |z|$ and $(-z)(-z)^\top = z z^\top$, we have $\Pi(-z) = \Pi(z)$ and $|{-z}|^{-1} = |z|^{-1}$, so $A(-z) = A(z)$.
\end{proof}

\begin{lemma}[Positive semidefiniteness of $A$]
\label{lem:A_psd}
For all $z \in \R^3 \setminus \{0\}$, the matrix $A(z)$ is positive semidefinite. Specifically, for any $Y \in \R^3$:
\[
Y^\top A(z) Y = \Psi(|z|)\big(|z|^2 |Y|^2 - (z \cdot Y)^2\big) \geq 0.
\]
Equality holds if and only if $Y$ is parallel to $z$ (i.e., $Y = \lambda z$ for some $\lambda \in \R$).
\end{lemma}

\begin{proof}
We compute:
\[
Y^\top A(z) Y = \Psi(|z|) \, Y^\top (|z|^2 I_3 - z z^\top) Y = \Psi(|z|)\big(|z|^2 |Y|^2 - (z \cdot Y)^2\big).
\]
By the Cauchy--Schwarz inequality, $(z \cdot Y)^2 \leq |z|^2 |Y|^2$, so $Y^\top A(z) Y \geq 0$. Equality in Cauchy--Schwarz holds if and only if $Y$ and $z$ are linearly dependent.
\end{proof}

\begin{lemma}[Projection annihilation: $z^\top A(z) = 0$]
\label{lem:zA_zero}
For all $z \in \R^3 \setminus \{0\}$:
\[
z^\top A(z) = 0^\top \qquad \text{and} \qquad A(z) \, z = 0.
\]
\end{lemma}

\begin{proof}
We compute $z^\top \Pi(z) = z^\top - z^\top \frac{z z^\top}{|z|^2} = z^\top - \frac{|z|^2}{|z|^2} z^\top = 0^\top$. Multiplying by $|z|^{-1}$ gives $z^\top A(z) = 0$. The second identity follows by transposing, using $A(z)^\top = A(z)$.
\end{proof}

%=============================================================================
\section{The H-Theorem and Nullspace of the Landau Operator}
\label{sec:htheorem}
%=============================================================================

This section contains the detailed proof that the nullspace of the Landau collision operator consists exactly of local Maxwellians. This is the content that bridges Step~1 of the main theorem.

\begin{lemma}[Symmetrized weak formulation of the Landau operator]
\label{lem:symmetrized_weak}
Let $f : \R^3 \to (0, \infty)$ be sufficiently smooth with sufficient decay. For any smooth test function $\phi : \R^3 \to \R$:
\[
\int_{\R^3} Q(f, f)(v) \, \phi(v) \, dv = -\frac{1}{2} \iint_{\R^3 \times \R^3} \big(\nabla \phi(v) - \nabla \phi(w)\big)^\top A(v - w) \big[f(w) \nabla f(v) - f(v) \nabla f(w)\big] \, dv \, dw.
\]
\end{lemma}

\begin{proof}
Starting from the definition of $Q(f,f)$, multiply by $\phi(v)$ and integrate over $v$. Integration by parts (using the divergence form and the decay of $f$) moves $\nabla_v$ onto $\phi$:
\[
\int_{\R^3} Q(f,f) \, \phi \, dv = -\iint_{\R^3 \times \R^3} \nabla_v \phi(v) \cdot A(v-w) \big[f(w) \nabla_v f(v) - f(v) \nabla_w f(w)\big] \, dw \, dv.
\]
Swap the dummy variables $v \leftrightarrow w$. Using $A(w-v) = A(v-w)$ (Lemma~\ref{lem:A_symmetric}(b)), and noting that the bracket $[f(w)\nabla_v f(v) - f(v)\nabla_w f(w)]$ changes sign under the swap, we obtain:
\[
\int_{\R^3} Q(f,f) \, \phi \, dv = +\iint_{\R^3 \times \R^3} \nabla_w \phi(w) \cdot A(v-w) \big[f(w) \nabla_v f(v) - f(v) \nabla_w f(w)\big] \, dv \, dw.
\]
Adding these two representations and dividing by 2 yields the symmetrized form.
\end{proof}

\begin{lemma}[Entropy dissipation formula]
\label{lem:entropy_dissipation}
Let $f : \R^3 \to (0, \infty)$ be sufficiently smooth with sufficient decay. Define the \emph{entropy dissipation functional}:
\[
D(f) = \int_{\R^3} Q(f, f)(v) \, \log f(v) \, dv.
\]
Then:
\[
D(f) = -\frac{1}{2} \iint_{\R^3 \times \R^3} f(v) f(w) \big(\nabla \log f(v) - \nabla \log f(w)\big)^\top A(v-w) \big(\nabla \log f(v) - \nabla \log f(w)\big) \, dv \, dw.
\]
\end{lemma}

\begin{proof}
Apply Lemma~\ref{lem:symmetrized_weak} with $\phi(v) = \log f(v)$, so $\nabla \phi(v) = \nabla \log f(v)$. In the bracket, use the identity $\nabla f = f \nabla \log f$ to factor:
\[
f(w) \nabla f(v) - f(v) \nabla f(w) = f(v) f(w) \big(\nabla \log f(v) - \nabla \log f(w)\big).
\]
Substituting into the symmetrized weak form gives the result.
\end{proof}

\begin{theorem}[H-theorem for the Landau operator]
\label{thm:H_theorem}
Let $f : \R^3 \to (0, \infty)$ be sufficiently smooth with sufficient decay. Then:
\[
D(f) = \int_{\R^3} Q(f, f)(v) \, \log f(v) \, dv \leq 0.
\]
\end{theorem}

\begin{proof}
By Lemma~\ref{lem:entropy_dissipation}, $D(f)$ is the negative of a double integral of the quadratic form $Y^\top A(z) Y$ with $Y = \nabla \log f(v) - \nabla \log f(w)$ and $z = v - w$, weighted by the strictly positive function $f(v) f(w)$. By Lemma~\ref{lem:A_psd}, $Y^\top A(z) Y \geq 0$, so the integrand is non-negative, and therefore $D(f) \leq 0$.
\end{proof}

\begin{lemma}[Characterization of $D(f) = 0$: the functional equation]
\label{lem:D_zero_functional_eq}
Let $f : \R^3 \to (0, \infty)$ be sufficiently smooth with sufficient decay. If $D(f) = 0$, then for all $v, w \in \R^3$ with $v \neq w$, the vector $\nabla \log f(v) - \nabla \log f(w)$ is parallel to $v - w$. That is, there exists a scalar function $\lambda : \R^3 \times \R^3 \to \R$ such that:
\[
\nabla \log f(v) - \nabla \log f(w) = \lambda(v, w)(v - w) \qquad \text{for all } v, w \in \R^3.
\]
\end{lemma}

\begin{proof}
By Lemma~\ref{lem:entropy_dissipation}, $D(f) = 0$ means the integral of the non-negative integrand $f(v)f(w) \cdot (\nabla \log f(v) - \nabla \log f(w))^\top A(v-w) (\nabla \log f(v) - \nabla \log f(w))$ vanishes. Since $f(v)f(w) > 0$ for all $v, w$, the quadratic form $Y^\top A(z) Y$ must vanish for all $v, w$ (with $v \neq w$). By the equality condition in Lemma~\ref{lem:A_psd}, $Y = \nabla \log f(v) - \nabla \log f(w)$ must be parallel to $z = v - w$.
\end{proof}

\begin{lemma}[Solution of the functional equation: $\nabla \log f$ is affine]
\label{lem:functional_eq_solution}
Let $g : \R^3 \to \R^3$ be a smooth vector field satisfying:
\[
g(v) - g(w) = \lambda(v, w)(v - w) \qquad \text{for all } v, w \in \R^3,
\]
for some scalar function $\lambda$. Then there exist a constant vector $b \in \R^3$ and a constant $c_0 \in \R$ such that:
\[
g(v) = b + 2c_0 \, v \qquad \text{for all } v \in \R^3.
\]
\end{lemma}

\begin{proof}
The proof proceeds in three steps.

\textbf{Step 1: Reduction using $w = 0$.}
Setting $w = 0$ gives $g(v) = g(0) + \lambda(v, 0) \, v$. Let $b = g(0)$ and $\tilde{\lambda}(v) = \lambda(v, 0)$. Then:
\[
g(v) = b + \tilde{\lambda}(v) \, v.
\]

\textbf{Step 2: Curl-free constraint forces spherical symmetry of $\tilde{\lambda}$.}
Suppose additionally that $g = \nabla \psi$ for some scalar function $\psi$ (i.e., $g$ is curl-free: $\nabla \times g = 0$). Taking the curl of $g(v) = b + \tilde{\lambda}(v) v$:
\[
\nabla \times g = \nabla \tilde{\lambda} \times v = 0.
\]
(Here we used $\nabla \times b = 0$ and $\nabla \times v = 0$, and the product rule $\nabla \times (\tilde{\lambda} v) = \nabla \tilde{\lambda} \times v + \tilde{\lambda} (\nabla \times v)$.)

The condition $\nabla \tilde{\lambda} \times v = 0$ for all $v \neq 0$ means $\nabla \tilde{\lambda}(v)$ is parallel to $v$ for all $v$. This implies $\tilde{\lambda}$ depends only on $|v|$: there exists a function $\hat{\lambda} : [0, \infty) \to \R$ such that $\tilde{\lambda}(v) = \hat{\lambda}(|v|^2)$. We write $\tilde{\lambda}(v) = 2c(|v|^2)$ for convenience.

Thus $g(v) = b + 2c(|v|^2) \, v$.

\textbf{Step 3: $c$ is constant.}
Substitute $g(v) = b + 2c(|v|^2) v$ into the original equation $g(v) - g(w) = \lambda(v,w)(v - w)$:
\[
2c(|v|^2) v - 2c(|w|^2) w = \lambda(v,w)(v - w).
\]
Rearranging:
\[
\big[2c(|v|^2) - \lambda(v,w)\big] v = \big[2c(|w|^2) - \lambda(v,w)\big] w.
\]
Choose $v, w$ that are linearly independent (e.g., $v = e_1$, $w = e_2$). Since a linear combination of linearly independent vectors can only vanish if both coefficients are zero:
\[
2c(|v|^2) = \lambda(v,w) = 2c(|w|^2).
\]
Since this holds for arbitrary magnitudes $|v|$ and $|w|$, the function $c$ must be a constant $c_0 \in \R$. Therefore $g(v) = b + 2c_0 v$.
\end{proof}

\begin{remark}
\label{rem:curl_free}
In the application to kinetic theory, $g = \nabla \log f$ is automatically a gradient field, so the curl-free condition in Step~2 holds by construction.
\end{remark}

\begin{lemma}[Integration: $\log f$ is a polynomial of degree $\leq 2$]
\label{lem:log_f_quadratic}
If $\nabla \log f(v) = b + 2c_0 v$ for constants $b \in \R^3$ and $c_0 \in \R$, then there exists a constant $a_0 \in \R$ such that:
\[
\log f(v) = a_0 + b \cdot v + c_0 |v|^2.
\]
Equivalently, $\log f$ lies in $\operatorname{span}\{1, v_1, v_2, v_3, |v|^2\}$.
\end{lemma}

\begin{proof}
Integrating $\partial_{v_i} \log f = b_i + 2c_0 v_i$ component-wise gives $\log f(v) = a_0 + \sum_i b_i v_i + c_0 \sum_i v_i^2 = a_0 + b \cdot v + c_0 |v|^2$.
\end{proof}

\begin{theorem}[Nullspace of the Landau operator -- necessity]
\label{thm:nullspace_necessity}
Let $f : \R^3 \to (0, \infty)$ be sufficiently smooth with sufficient decay and finite entropy. If $Q(f, f) = 0$, then $\log f \in \operatorname{span}\{1, v_1, v_2, v_3, |v|^2\}$. That is, there exist constants $a_0 \in \R$, $b \in \R^3$, and $c_0 \in \R$ such that:
\[
\log f(v) = a_0 + b \cdot v + c_0 |v|^2.
\]
If additionally $f \in L^1(\R^3)$ (finite total mass), then $c_0 < 0$, and $f$ is a Maxwellian:
\[
f(v) = \frac{n}{(2\pi T)^{3/2}} \exp\!\left(-\frac{|v - u|^2}{2T}\right)
\]
for some $n > 0$, $u \in \R^3$, and $T > 0$ determined by $a_0, b, c_0$.
\end{theorem}

\begin{proof}
If $Q(f,f) = 0$, then for any test function $\phi$, $\int Q(f,f) \phi \, dv = 0$. In particular, choosing $\phi = \log f$: $D(f) = 0$. By Lemma~\ref{lem:D_zero_functional_eq}, $\nabla \log f(v) - \nabla \log f(w) = \lambda(v,w)(v-w)$. Since $\nabla \log f$ is a gradient field (Remark~\ref{rem:curl_free}), Lemma~\ref{lem:functional_eq_solution} gives $\nabla \log f(v) = b + 2c_0 v$. Lemma~\ref{lem:log_f_quadratic} yields $\log f(v) = a_0 + b \cdot v + c_0 |v|^2$.

For $f \in L^1(\R^3)$, the integral $\int_{\R^3} \exp(a_0 + b \cdot v + c_0 |v|^2) \, dv$ must converge, which requires $c_0 < 0$. Completing the square: $a_0 + b \cdot v + c_0 |v|^2 = (a_0 + |b|^2/(4|c_0|)) + c_0 |v - u|^2$ with $u = -b/(2c_0)$. Setting $T = -1/(2c_0)$ and $n = (2\pi T)^{3/2} \exp(a_0 + |b|^2/(4|c_0|))$ gives the Maxwellian form.
\end{proof}

\begin{theorem}[Nullspace of the Landau operator -- sufficiency]
\label{thm:nullspace_sufficiency}
If $\log f(v) = a_0 + b \cdot v + c_0 |v|^2$ for constants $a_0 \in \R$, $b \in \R^3$, $c_0 \in \R$ (with $c_0 < 0$ for integrability), then $Q(f, f) = 0$.
\end{theorem}

\begin{proof}
We have $\nabla \log f(v) = b + 2c_0 v$. The relative gradient is:
\[
\nabla \log f(v) - \nabla \log f(w) = 2c_0(v - w).
\]
In the score form of the Landau operator:
\[
Q(f,f)(v) = \nabla_v \cdot \left\{ f(v) \int_{\R^3} A(v-w) f(w) \cdot 2c_0(v-w) \, dw \right\}.
\]
By Lemma~\ref{lem:zA_zero}, $A(v-w)(v-w) = 0$ for all $v \neq w$. Therefore the integrand vanishes identically, and $Q(f,f) = 0$.
\end{proof}

\begin{corollary}[Complete characterization of the nullspace]
\label{cor:nullspace}
For $f : \R^3 \to (0, \infty)$ sufficiently smooth with sufficient decay and $f \in L^1(\R^3)$:
\[
Q(f, f) = 0 \quad \iff \quad f \text{ is a Maxwellian distribution.}
\]
\end{corollary}

\begin{proof}
Combine Theorem~\ref{thm:nullspace_necessity} and Theorem~\ref{thm:nullspace_sufficiency}.
\end{proof}

%=============================================================================
\section{Vlasov--Maxwell Transport Constraints}
\label{sec:transport}
%=============================================================================

This section shows how the Vlasov--Maxwell transport equations, combined with the local Maxwellian form from Section~\ref{sec:htheorem}, force a global (spatially uniform) Maxwellian.

\begin{lemma}[Vanishing of the left-hand side of the entropy identity at steady state]
\label{lem:lhs_vanishes}
Let $(f, E, B)$ be a steady state of the VML system (Definition~\ref{def:steady_state}) on $\T^3 \times \R^3$. Then:
\[
\iint_{\T^3 \times \R^3} \big[v \cdot \nabla_x f + (E + v \times B) \cdot \nabla_v f\big] \log f \, dv \, dx = 0.
\]
\end{lemma}

\begin{proof}
Three terms must be shown to vanish.

\textbf{Spatial transport:} Using $\log f \, \nabla_x f = \nabla_x(f \log f - f)$,
\[
\iint v \cdot \nabla_x f \, \log f \, dv \, dx = \iint v \cdot \nabla_x(f \log f - f) \, dv \, dx = \int_{\R^3} v \cdot \underbrace{\int_{\T^3} \nabla_x(f \log f - f) \, dx}_{= \, 0 \text{ by periodicity}} \, dv = 0.
\]

\textbf{Electric force:} Integrating by parts in $v$,
\[
\iint E \cdot \nabla_v f \, \log f \, dv \, dx = -\iint E \cdot f \, \nabla_v \log f \, dv \, dx + \text{b.t.} = -\iint f \, \underbrace{\nabla_v \cdot E}_{= \, 0} \, dv \, dx = 0.
\]
(We used $\nabla_v f \, \log f = \nabla_v(f \log f - f)$ and integrated by parts; or equivalently, $E \cdot \nabla_v f \, \log f$ integrated by parts gives $-f \, \nabla_v \cdot E = 0$ since $E = E(x)$ does not depend on $v$.)

\textbf{Magnetic force:} Similarly,
\[
\iint (v \times B) \cdot \nabla_v f \, \log f \, dv \, dx = -\iint f \, \underbrace{\nabla_v \cdot (v \times B)}_{= \, 0} \, dv \, dx = 0.
\]
Here $\nabla_v \cdot (v \times B) = B \cdot (\nabla_v \times v) - v \cdot (\nabla_v \times B)$. Since $\nabla_v \times v = 0$ (the curl of the identity vector field is zero in 3D; more precisely, $(\nabla_v \times v)_i = \sum_{jk} \epsilon_{ijk} \partial_{v_j} v_k = \sum_{jk} \epsilon_{ijk} \delta_{jk} = \sum_j \epsilon_{ijj} = 0$) and $B$ does not depend on $v$, this vanishes.
\end{proof}

\begin{lemma}[Global entropy production vanishes at steady state]
\label{lem:global_entropy_zero}
Let $(f, E, B)$ be a steady state of the VML system with $\nu > 0$. Then:
\[
\int_{\T^3} D_x(f) \, dx = 0,
\]
where $D_x(f) = \int_{\R^3} Q(f,f)(x, \cdot)(v) \log f(x, v) \, dv \leq 0$ is the local entropy dissipation at spatial point $x$.
\end{lemma}

\begin{proof}
At steady state, the VML equation gives $v \cdot \nabla_x f + (E + v \times B) \cdot \nabla_v f = \nu Q[f,f]$. Multiplying by $\log f$ and integrating over $\T^3 \times \R^3$:
\[
\underbrace{\iint [\text{l.h.s.}] \log f \, dv \, dx}_{= \, 0 \text{ (Lemma~\ref{lem:lhs_vanishes})}} = \nu \iint Q[f,f] \log f \, dv \, dx = \nu \int_{\T^3} D_x(f) \, dx.
\]
Since $\nu > 0$, we get $\int_{\T^3} D_x(f) \, dx = 0$.
\end{proof}

\begin{lemma}[Pointwise vanishing of entropy dissipation]
\label{lem:pointwise_D_zero}
If $D_x(f) \leq 0$ for all $x \in \T^3$ (Theorem~\ref{thm:H_theorem}) and $\int_{\T^3} D_x(f) \, dx = 0$ (Lemma~\ref{lem:global_entropy_zero}), then $D_x(f) = 0$ for all $x \in \T^3$.
\end{lemma}

\begin{proof}
A non-positive continuous function whose integral over a domain of positive measure vanishes must be identically zero.
\end{proof}

\begin{corollary}[Steady state is a local Maxwellian]
\label{cor:local_maxwellian}
At any steady state of the VML system with $\nu > 0$, the distribution $f(x, v)$ is a local Maxwellian (Definition~\ref{def:local_maxwellian}) for each $x \in \T^3$.
\end{corollary}

\begin{proof}
By Lemma~\ref{lem:pointwise_D_zero}, $D_x(f) = 0$ for all $x$. By Corollary~\ref{cor:nullspace}, $f(x, \cdot)$ is a Maxwellian for each $x$.
\end{proof}

%=============================================================================
\section{Polynomial Matching: From Local to Global Maxwellian}
\label{sec:polynomial}
%=============================================================================

Since $Q[f,f] = 0$ at steady state, the Vlasov equation reduces to the collisionless transport:
\[
v \cdot \nabla_x(\log f) + (E + v \times B) \cdot \nabla_v(\log f) = 0.
\]

\begin{lemma}[Polynomial identity in velocity]
\label{lem:polynomial_identity}
Let $f$ be a local Maxwellian with $\log f(x,v) = a(x) + b(x) \cdot v + c(x)|v|^2$ (where $a, b, c$ are related to $n, u, T$ as in Definition~\ref{def:local_maxwellian}; specifically $c(x) = -1/(2T(x))$ and $b(x) = u(x)/T(x)$). At steady state, the collisionless transport equation is equivalent to:
\begin{equation}
\label{eq:polynomial}
\underbrace{(v \cdot \nabla_x c)|v|^2}_{\mathcal{O}(|v|^3)} + \underbrace{\sum_{i,j=1}^{3} v_i v_j \, \partial_{x_i} b_j}_{\mathcal{O}(|v|^2)} + \underbrace{v \cdot \big(\nabla_x a + 2c \, E + B \times b\big)}_{\mathcal{O}(|v|^1)} + \underbrace{E \cdot b}_{\mathcal{O}(|v|^0)} = 0.
\end{equation}
This must hold for all $v \in \R^3$.
\end{lemma}

\begin{proof}
Compute $\nabla_x \log f = \nabla_x a + (\nabla_x b) v + (\nabla_x c)|v|^2$ (where $\nabla_x b$ is the Jacobian matrix $(\partial_{x_i} b_j)_{ij}$) and $\nabla_v \log f = b + 2cv$. Substituting:
\begin{align*}
v \cdot \nabla_x \log f &= v \cdot \nabla_x a + \sum_{i,j} v_i v_j \partial_{x_i} b_j + (v \cdot \nabla_x c)|v|^2, \\
E \cdot \nabla_v \log f &= E \cdot b + 2c \, E \cdot v, \\
(v \times B) \cdot \nabla_v \log f &= (v \times B) \cdot b + 2c \underbrace{(v \times B) \cdot v}_{= \, 0}.
\end{align*}
Using $(v \times B) \cdot b = v \cdot (B \times b)$ (cyclic property of the scalar triple product), and collecting by degree in $v$, we obtain \eqref{eq:polynomial}.
\end{proof}

\begin{lemma}[Temperature is spatially constant]
\label{lem:T_constant}
Under the conditions of Lemma~\ref{lem:polynomial_identity}, $\nabla_x c = 0$, i.e., $T(x) \equiv T_\infty$ is a global constant.
\end{lemma}

\begin{proof}
The $\mathcal{O}(|v|^3)$ terms in \eqref{eq:polynomial} give $(v \cdot \nabla_x c)|v|^2 = 0$ for all $v$. Choosing $v = t \, e_i$ for $t \to \infty$ shows $\partial_{x_i} c = 0$ for each $i$. Since $c = -1/(2T)$, $T$ is constant.
\end{proof}

\begin{lemma}[Bulk velocity is spatially constant]
\label{lem:u_constant}
Under the conditions of Lemma~\ref{lem:polynomial_identity} with $c$ constant, $\partial_{x_i} b_j = 0$ for all $i, j$, i.e., $b(x) \equiv b_\infty$ and $u(x) \equiv u_\infty$ are global constants.
\end{lemma}

\begin{proof}
With $\nabla_x c = 0$, the $\mathcal{O}(|v|^2)$ terms give $\sum_{i,j} v_i v_j \partial_{x_i} b_j = 0$ for all $v$. Setting $v = e_k$: $\partial_{x_k} b_k = 0$ (no sum). Setting $v = e_k + e_l$ with $k \neq l$: $\partial_{x_k} b_k + \partial_{x_k} b_l + \partial_{x_l} b_k + \partial_{x_l} b_l = 0$, which gives $\partial_{x_k} b_l + \partial_{x_l} b_k = 0$ (Killing's equation).

From Killing's equation on the flat torus: $\partial_{x_k} b_l = -\partial_{x_l} b_k$. Taking $\partial_{x_k}$: $\partial_{x_k}^2 b_l = -\partial_{x_k} \partial_{x_l} b_k$. Taking $\partial_{x_l}$ of $\partial_{x_l} b_k = -\partial_{x_k} b_l$: $\partial_{x_l}^2 b_k = -\partial_{x_l} \partial_{x_k} b_l$. So $\partial_{x_k}^2 b_l = \partial_{x_l}^2 b_k$. Summing over $k$: $\Delta b_l = \partial_{x_l}(\nabla_x \cdot b)$. But $\nabla_x \cdot b = \sum_k \partial_{x_k} b_k = 0$ (from the diagonal terms). So $\Delta b_l = 0$ for each $l$.

On the flat torus, integrate: $\int_{\T^3} |\nabla_x b_l|^2 \, dx = -\int_{\T^3} b_l \Delta b_l \, dx = 0$ (by parts, using periodicity). Therefore $\nabla_x b_l = 0$ for each $l$, and $b$ is constant.
\end{proof}

\begin{lemma}[Force balance equation]
\label{lem:force_balance}
Under the conditions of Lemmas~\ref{lem:T_constant} and \ref{lem:u_constant} (constant $T_\infty$, constant $u_\infty$), the $\mathcal{O}(|v|^1)$ terms of \eqref{eq:polynomial} yield:
\[
\nabla_x \log n(x) = \frac{1}{T_\infty}\big(E(x) + u_\infty \times B(x)\big).
\]
\end{lemma}

\begin{proof}
With $c$ and $b$ constant, the $\mathcal{O}(|v|^1)$ coefficient in \eqref{eq:polynomial} is $\nabla_x a + 2c E + B \times b = 0$. Since $a(x) = \log n(x) - \frac{3}{2}\log(2\pi T_\infty) + |u_\infty|^2/(2T_\infty)$ (from the Maxwellian form), we have $\nabla_x a = \nabla_x \log n$. Substituting $c = -1/(2T_\infty)$ and $b = u_\infty/T_\infty$:
\[
\nabla_x \log n - \frac{E}{T_\infty} + B \times \frac{u_\infty}{T_\infty} = 0.
\]
Using $B \times u_\infty = -u_\infty \times B$, rearranging gives the result.
\end{proof}

\begin{lemma}[Zeroth-order term]
\label{lem:zeroth_order}
The $\mathcal{O}(|v|^0)$ term of \eqref{eq:polynomial} gives $E(x) \cdot b_\infty = 0$, i.e.:
\[
E(x) \cdot u_\infty = 0 \qquad \text{for all } x \in \T^3.
\]
\end{lemma}

\begin{proof}
Direct reading of the constant term in \eqref{eq:polynomial}, using $b_\infty = u_\infty / T_\infty$ and $T_\infty > 0$.
\end{proof}

%=============================================================================
\section{Nullification of Bulk Velocity via Amp\`ere's Law}
\label{sec:ampere}
%=============================================================================

\begin{lemma}[Amp\`ere's law at steady state]
\label{lem:ampere_steady}
At steady state ($\partial_t E = 0$), Amp\`ere's law \eqref{eq:ampere} gives:
\[
\nabla_x \times B = J = n(x) \, u_\infty.
\]
\end{lemma}

\begin{proof}
At steady state $\partial_t E = 0$, so \eqref{eq:ampere} becomes $\nabla_x \times B = J$. For a local Maxwellian with constant bulk velocity $u_\infty$, $J(x) = \int v f \, dv = n(x) u_\infty$.
\end{proof}

\begin{lemma}[Vanishing of bulk velocity]
\label{lem:u_zero}
Under the conditions of Lemma~\ref{lem:ampere_steady}, with $n(x) > 0$ for all $x \in \T^3$:
\[
u_\infty = 0.
\]
\end{lemma}

\begin{proof}
Take the dot product of $\nabla_x \times B = n(x) u_\infty$ with $u_\infty$:
\[
u_\infty \cdot (\nabla_x \times B) = n(x) |u_\infty|^2.
\]
Using the vector identity $\nabla_x \cdot (B \times u_\infty) = u_\infty \cdot (\nabla_x \times B) - B \cdot (\nabla_x \times u_\infty)$ and $\nabla_x \times u_\infty = 0$ (since $u_\infty$ is constant):
\[
\nabla_x \cdot (B \times u_\infty) = n(x)|u_\infty|^2.
\]
Integrate over the periodic domain $\T^3$. The left-hand side vanishes by the divergence theorem on a closed manifold:
\[
0 = |u_\infty|^2 \int_{\T^3} n(x) \, dx.
\]
Since $n(x) > 0$, the integral $\int_{\T^3} n(x) \, dx > 0$. Therefore $|u_\infty|^2 = 0$, i.e., $u_\infty = 0$.
\end{proof}

%=============================================================================
\section{Spatial Uniformity via the Maximum Principle}
\label{sec:maximum_principle}
%=============================================================================

\begin{lemma}[Poisson--Boltzmann equation for the density]
\label{lem:poisson_boltzmann}
With $u_\infty = 0$, the force balance (Lemma~\ref{lem:force_balance}) and Gauss's law \eqref{eq:gauss} yield:
\[
T_\infty \, \Delta_x(\log n) = n(x) - \rho_{\mathrm{ion}}.
\]
\end{lemma}

\begin{proof}
With $u_\infty = 0$, the force balance gives $\nabla_x \log n = E / T_\infty$. Taking the divergence: $\Delta_x \log n = (\nabla_x \cdot E)/T_\infty$. By Gauss's law, $\nabla_x \cdot E = n - \rho_{\mathrm{ion}}$. The result follows.
\end{proof}

\begin{lemma}[Maximum principle: density is constant]
\label{lem:density_constant}
Let $n : \T^3 \to (0, \infty)$ be smooth and satisfy $T_\infty \Delta_x(\log n) = n - \rho_{\mathrm{ion}}$ with $T_\infty > 0$ and $\rho_{\mathrm{ion}} > 0$. Then $n(x) \equiv \rho_{\mathrm{ion}}$.
\end{lemma}

\begin{proof}
Since $\T^3$ is compact and $n$ is smooth and positive, $n$ achieves a global maximum at some $x_{\max}$ and a global minimum at some $x_{\min}$.

At $x_{\max}$: $\log n$ also achieves its maximum (since $\log$ is monotone increasing), so $\Delta_x(\log n)(x_{\max}) \leq 0$. Therefore $n(x_{\max}) - \rho_{\mathrm{ion}} = T_\infty \Delta_x(\log n)(x_{\max}) \leq 0$, giving $n(x_{\max}) \leq \rho_{\mathrm{ion}}$.

At $x_{\min}$: similarly, $\Delta_x(\log n)(x_{\min}) \geq 0$, so $n(x_{\min}) \geq \rho_{\mathrm{ion}}$.

Combining: $\rho_{\mathrm{ion}} \leq n(x_{\min}) \leq n(x) \leq n(x_{\max}) \leq \rho_{\mathrm{ion}}$ for all $x$, so $n(x) \equiv \rho_{\mathrm{ion}}$.
\end{proof}

\begin{corollary}[Electric field vanishes]
\label{cor:E_zero}
Under the conditions of Lemma~\ref{lem:density_constant}, $E(x) = 0$ for all $x \in \T^3$.
\end{corollary}

\begin{proof}
Since $n \equiv \rho_{\mathrm{ion}}$, we have $\nabla_x \log n = 0$. The force balance with $u_\infty = 0$ gives $E = T_\infty \nabla_x \log n = 0$.
\end{proof}

%=============================================================================
\section{Uniformity of the Magnetic Field}
\label{sec:magnetic}
%=============================================================================

\begin{lemma}[Magnetic field is spatially constant]
\label{lem:B_constant}
With $u_\infty = 0$ and $J = n u_\infty = 0$, Amp\`ere's law gives $\nabla_x \times B = 0$. Combined with $\nabla_x \cdot B = 0$, this implies $B(x) \equiv B_\infty$ is constant on $\T^3$.
\end{lemma}

\begin{proof}
From $\nabla_x \times B = 0$ and $\nabla_x \cdot B = 0$, the vector Laplacian is:
\[
\Delta_x B = \nabla_x(\nabla_x \cdot B) - \nabla_x \times (\nabla_x \times B) = 0 - 0 = 0.
\]
So each component $B_i$ is harmonic on $\T^3$. Multiplying $\Delta B_i = 0$ by $B_i$ and integrating by parts on $\T^3$ (using periodicity):
\[
0 = \int_{\T^3} B_i \, \Delta B_i \, dx = -\int_{\T^3} |\nabla_x B_i|^2 \, dx.
\]
Therefore $\nabla_x B_i = 0$ for each $i$, so $B$ is constant.
\end{proof}

\begin{lemma}[Uniform magnetic field is compatible with Maxwellian equilibrium]
\label{lem:B_compatible}
For any constant $B_\infty \in \R^3$ and any isotropic Maxwellian $f_\infty(v) \propto \exp(-|v|^2/(2T_\infty))$:
\[
(v \times B_\infty) \cdot \nabla_v f_\infty = 0 \qquad \text{for all } v \in \R^3.
\]
\end{lemma}

\begin{proof}
$\nabla_v f_\infty = -\frac{v}{T_\infty} f_\infty$, so $(v \times B_\infty) \cdot \nabla_v f_\infty = -\frac{1}{T_\infty} (v \times B_\infty) \cdot v \, f_\infty = 0$, since $v \times B_\infty$ is orthogonal to $v$.
\end{proof}

%=============================================================================
\section{Conservation Laws and Determination of Equilibrium Parameters}
\label{sec:conservation}
%=============================================================================

\begin{lemma}[Conservation of spatial mean of $B$]
\label{lem:B_mean_conserved}
Under Faraday's law $\partial_t B = -\nabla_x \times E$ on the periodic domain $\T^3$:
\[
\frac{d}{dt} \int_{\T^3} B(t, x) \, dx = 0.
\]
\end{lemma}

\begin{proof}
$\frac{d}{dt} \int_{\T^3} B \, dx = -\int_{\T^3} \nabla_x \times E \, dx = 0$ by Stokes' theorem on the closed manifold $\T^3$ (equivalently, by periodicity).
\end{proof}

\begin{lemma}[Conservation of total energy]
\label{lem:energy_conserved}
The total energy $\mathcal{E}(t) = \iint \frac{1}{2}|v|^2 f \, dv \, dx + \frac{1}{2}\int(|E|^2 + |B|^2) \, dx$ is conserved under the VML dynamics.
\end{lemma}

\begin{proof}[Proof sketch]
Multiply the Vlasov equation by $\frac{1}{2}|v|^2$ and integrate. The collision term contributes zero (since $|v|^2$ is a collisional invariant of the Landau operator). The transport term integrates to zero by periodicity. The Lorentz force term gives $\frac{d}{dt}(\text{kinetic energy}) = \int J \cdot E \, dx$ (since $v \cdot (v \times B) = 0$). From Maxwell's equations: $\frac{d}{dt}\frac{1}{2}\int(|E|^2 + |B|^2) dx = -\int J \cdot E \, dx$. Adding the two gives $\frac{d\mathcal{E}}{dt} = 0$.
\end{proof}

\begin{lemma}[Determination of $B_\infty$]
\label{lem:B_infty}
At equilibrium, $B(x) \equiv B_\infty$ is constant. By Lemma~\ref{lem:B_mean_conserved}:
\[
B_\infty = \frac{1}{|\T^3|} \int_{\T^3} B_{\mathrm{init}}(x) \, dx.
\]
\end{lemma}

\begin{lemma}[Determination of $T_\infty$]
\label{lem:T_infty}
At equilibrium, $f_\infty(v) = \frac{\rho_{\mathrm{ion}}}{(2\pi T_\infty)^{3/2}} \exp\!\big(-\frac{|v|^2}{2T_\infty}\big)$, $E = 0$, and $B \equiv B_\infty$. The total energy is:
\[
\mathcal{E} = \frac{3}{2} \rho_{\mathrm{ion}} \, T_\infty \, |\T^3| + \frac{|B_\infty|^2}{2} |\T^3|.
\]
By conservation (Lemma~\ref{lem:energy_conserved}), $\mathcal{E} = \mathcal{E}_0$. Solving:
\[
T_\infty = \frac{2}{3 \, \rho_{\mathrm{ion}} \, |\T^3|} \left(\mathcal{E}_0 - \frac{|B_\infty|^2}{2}|\T^3|\right).
\]
\end{lemma}

\begin{lemma}[Positivity of $T_\infty$]
\label{lem:T_positive}
$T_\infty > 0$, which is equivalent to $\mathcal{E}_0 > \frac{|B_\infty|^2}{2}|\T^3|$. This holds by Jensen's inequality:
\[
|B_\infty|^2 = \left|\frac{1}{|\T^3|}\int_{\T^3} B_{\mathrm{init}} \, dx\right|^2 \leq \frac{1}{|\T^3|} \int_{\T^3} |B_{\mathrm{init}}|^2 \, dx,
\]
so $\frac{|B_\infty|^2}{2}|\T^3| \leq \frac{1}{2}\int_{\T^3}|B_{\mathrm{init}}|^2 \, dx \leq \frac{1}{2}\int_{\T^3}(|E_{\mathrm{init}}|^2 + |B_{\mathrm{init}}|^2) \, dx < \mathcal{E}_0$ (the last strict inequality uses the fact that $\iint \frac{1}{2}|v|^2 f_{\mathrm{init}} \, dv \, dx > 0$ since $f_{\mathrm{init}} > 0$).
\end{lemma}

%=============================================================================
\section{Main Theorem: Assembly}
\label{sec:main}
%=============================================================================

\begin{theorem}[Global steady state of the VML system]
\label{thm:main}
Consider the Vlasov--Maxwell--Landau system (Definition~\ref{def:vml}) with:
\begin{itemize}
\item collision frequency $\nu > 0$,
\item periodic spatial domain $x \in \T^3$, velocity space $v \in \R^3$,
\item uniform neutralizing ion background density $\rho_{\mathrm{ion}} > 0$,
\item initial data $(f_{\mathrm{init}}, E_{\mathrm{init}}, B_{\mathrm{init}})$ with $f_{\mathrm{init}} > 0$, finite total energy $\mathcal{E}_0 > 0$, and finite entropy.
\end{itemize}
Let $(f, E, B)$ be a sufficiently smooth steady-state solution (Definition~\ref{def:steady_state}). Then:
\begin{enumerate}[label=(\roman*)]
\item The distribution is a spatially uniform, zero-drift Maxwellian:
\[
f_\infty(v) = \frac{\rho_{\mathrm{ion}}}{(2\pi T_\infty)^{3/2}} \exp\!\left(-\frac{|v|^2}{2T_\infty}\right).
\]

\item The electric field vanishes: $E_\infty = 0$.

\item The magnetic field is a spatial constant equal to the mean of the initial field:
\[
B_\infty = \frac{1}{|\T^3|}\int_{\T^3} B_{\mathrm{init}}(x) \, dx.
\]

\item The equilibrium temperature is uniquely determined:
\[
T_\infty = \frac{2}{3\, \rho_{\mathrm{ion}} \, |\T^3|} \left(\mathcal{E}_0 - \frac{|B_\infty|^2}{2} |\T^3|\right) > 0,
\]
where $\mathcal{E}_0 = \iint \frac{1}{2}|v|^2 f_{\mathrm{init}} \, dv \, dx + \frac{1}{2}\int (|E_{\mathrm{init}}|^2 + |B_{\mathrm{init}}|^2) \, dx$.
\end{enumerate}
\end{theorem}

\begin{proof}
The proof assembles the lemmas from the preceding sections:
\begin{enumerate}
\item \textbf{$f$ is a local Maxwellian:} Corollary~\ref{cor:local_maxwellian} (via Lemmas~\ref{lem:lhs_vanishes}, \ref{lem:global_entropy_zero}, \ref{lem:pointwise_D_zero}, and the H-theorem results of Section~\ref{sec:htheorem}).

\item \textbf{Temperature is constant:} Lemma~\ref{lem:T_constant} (from the $\mathcal{O}(|v|^3)$ terms).

\item \textbf{Bulk velocity is constant:} Lemma~\ref{lem:u_constant} (from the $\mathcal{O}(|v|^2)$ terms and Killing's equation).

\item \textbf{Bulk velocity is zero:} Lemma~\ref{lem:u_zero} (from Amp\`ere's law on the periodic domain).

\item \textbf{Density is constant and $E = 0$:} Lemmas~\ref{lem:poisson_boltzmann}, \ref{lem:density_constant}, Corollary~\ref{cor:E_zero} (from the force balance and maximum principle).

\item \textbf{Magnetic field is constant:} Lemma~\ref{lem:B_constant} (from $\nabla \times B = 0$ and $\nabla \cdot B = 0$).

\item \textbf{Parameters determined:} Lemmas~\ref{lem:B_infty}, \ref{lem:T_infty}, \ref{lem:T_positive} (from conservation laws).
\end{enumerate}
\end{proof}

%=============================================================================
\section{Dependency Graph}
\label{sec:dependencies}
%=============================================================================

The logical dependencies between the lemmas are:

\bigskip
\noindent
\textbf{Algebraic core} (no dependencies):
\begin{itemize}
\item Lemma~\ref{lem:A_symmetric} (symmetry of $A$)
\item Lemma~\ref{lem:A_psd} (positive semidefiniteness, uses Cauchy--Schwarz)
\item Lemma~\ref{lem:zA_zero} (projection annihilation)
\end{itemize}

\medskip
\noindent
\textbf{H-theorem chain} (sequential):
\[
\text{Lem.~\ref{lem:A_symmetric}} \to \text{Lem.~\ref{lem:symmetrized_weak}} \to \text{Lem.~\ref{lem:entropy_dissipation}} \to \text{Thm.~\ref{thm:H_theorem} (uses Lem.~\ref{lem:A_psd})}
\]

\medskip
\noindent
\textbf{Nullspace chain} (sequential):
\[
\text{Thm.~\ref{thm:H_theorem}} + \text{Lem.~\ref{lem:A_psd}} \to \text{Lem.~\ref{lem:D_zero_functional_eq}} \to \text{Lem.~\ref{lem:functional_eq_solution}} \to \text{Lem.~\ref{lem:log_f_quadratic}} \to \text{Thm.~\ref{thm:nullspace_necessity}}
\]
\[
\text{Lem.~\ref{lem:zA_zero}} \to \text{Thm.~\ref{thm:nullspace_sufficiency}}
\]

\medskip
\noindent
\textbf{Transport chain} (sequential):
\[
\text{Lem.~\ref{lem:lhs_vanishes}} \to \text{Lem.~\ref{lem:global_entropy_zero}} \to \text{Lem.~\ref{lem:pointwise_D_zero}} \to \text{Cor.~\ref{cor:local_maxwellian} (uses Cor.~\ref{cor:nullspace})}
\]

\medskip
\noindent
\textbf{Polynomial matching chain} (sequential, uses Cor.~\ref{cor:local_maxwellian}):
\[
\text{Lem.~\ref{lem:polynomial_identity}} \to \text{Lem.~\ref{lem:T_constant}} \to \text{Lem.~\ref{lem:u_constant}} \to \text{Lem.~\ref{lem:force_balance}} + \text{Lem.~\ref{lem:zeroth_order}}
\]

\medskip
\noindent
\textbf{Amp\`ere chain} (uses Lem.~\ref{lem:u_constant}):
\[
\text{Lem.~\ref{lem:ampere_steady}} \to \text{Lem.~\ref{lem:u_zero}}
\]

\medskip
\noindent
\textbf{Maximum principle chain} (uses Lem.~\ref{lem:u_zero} + Lem.~\ref{lem:force_balance}):
\[
\text{Lem.~\ref{lem:poisson_boltzmann}} \to \text{Lem.~\ref{lem:density_constant}} \to \text{Cor.~\ref{cor:E_zero}}
\]

\medskip
\noindent
\textbf{Magnetic + parameters} (uses Lem.~\ref{lem:u_zero}):
\[
\text{Lem.~\ref{lem:B_constant}} \to \text{Lem.~\ref{lem:B_infty}} + \text{Lem.~\ref{lem:T_infty}} + \text{Lem.~\ref{lem:T_positive}}
\]

\end{document}
